% !TEX root = ../main.tex
\frontchapter{Acknowledgments}

% \vspace{-1em}

% {\normalsize
%   \noindent \emph{I thank my thesis committee --- George Amvrosiadis (chair), Greg Ganger, Andy
%     Pavlo, Vyas Sekar, and Brad Settlemyer --- for their guidance and feedback throughout this work.
    % {\small \noindent \emph{I thank my thesis committee --- George Amvrosiadis (chair), Greg Ganger,
    % Andy Pavlo, Vyas Sekar, and Brad Settlemyer --- for their guidance and feedback throughout this
    % work. This research was supported by Los Alamos National Laboratory under contracts
    % [LANL-CONTRACT-1] and [LANL-CONTRACT-2], and by the U.S. Department of Energy under grants
    % [DOE-GRANT-1] and [DOE-GRANT-2]. I also thank the current and former members of the PDL Consortium
    % (Bloomberg LP, Datadog, Everpure, Google, Jane Street, LayerZero Labs, Meta, Microsoft Research,
    % Oracle Corporation, Salesforce, Uber, and Western Digital) for their interest, insights, and
    % feedback.} \par }
    % This research was sponsored by the Los Alamos National Laboratory's Institute for Reliable
    % High Performance Information Technology (IRHPIT), the NetSketch project, and the Mochi project
    % (DOE/NNSA Contract Nos.~DE-AC52-06NA25396/394903, DE-AC52-06NA25396/394957, and
    % 89233218CNA000001/394957 respectively), and the New Mexico Consortium.
    % I also thank the
    % current and former members of the Parallel Data Lab (PDL) Consortium at Carnegie Mellon
    % University (Bloomberg LP, Datadog, Everpure, Google, Jane Street, LayerZero Labs, Meta,
    % Microsoft Research, Oracle Corporation, Salesforce, Uber, and Western Digital) for their
    % interest, insights, and feedback.
    % The views and the conclusions are those of the author and do not necessarily represent the opinions
    % of the sponsoring organizations, agencies, or the U.S.
    % Government.
%   }
%   \par
% }

% {\normalsize
%   \noindent \emph{This thesis was supervised by a committee comprising George Amvrosiadis
%     (chair), Greg Ganger, Andy Pavlo, Vyas Sekar, and Brad Settlemyer. I thank them for their
%     guidance and feedback throughout this work.}
%   \par
% }

{\normalsize
  \noindent \emph{I thank my thesis committee --- George Amvrosiadis (chair), Greg Ganger, Andy
    Pavlo, Vyas Sekar, and Brad Settlemyer --- for their guidance and feedback throughout this
    work.}
  \par
}

\vspace{1em}

\looseness=-1
While the bulk of this
document is devoted to the technical contributions of
a scenic route through grad school, it is pertinent to take
a moment to acknowledge the substantial human element enabling this work.
The story arc necessitates that it begin with Saikat, my mentor at MSR
Bangalore from 2016--18, who prevailed in an argument about me not wanting to go to grad school.
While not entirely convinced, I deferred to his superior intellect and applied and surprisingly
got some admits.
Next, George, my advisor for the last 8 years.
The time and space he gave me to go after
misbehaving systems appeared questionable in the moment, but
proved absolutely decisive in enabling this work.
Chuck for being a wizard of a consultant for when anything caught fire: build systems,
fabrics and drivers, my code, other people's code, but somehow never his own.
Brad for being engaged throughout: it is rare and valued.
Qing for onboarding me with DeltaFS and remaining a close collaborator throughout.
Phil and Rob for both being very wise and for willing to make it accessible to me when needed.
Vyas for being a black box whose office one leaves with less uncertainty
than they enter with --- some forced add-ons in the form of lessons from Heilmeier and Walt
Disney not diminishing the overall deal value.

Next, the Parallel Data Lab --- the people, compute, and the calendar.
Greg, now tragically deposed and exiled to an obscure hyperscaler (Google MTV),
for an excellent job running the show for most of my tenure here.
Bill, Karen, and Joan for all the magical things they do to make everything run,
and Joan again for the last-minute typesetting lessons.
% Bill for always bolstering the audience numbers during practice talks, and him, Karen, Joan for
% all the magical things they do to make everything run.
Jason, Chad, and Mitch for handling all my tantrums when Wolf did not work,
and for all the times it did.
My appetite for the scale of systems I wanted to do
went down by a couple orders of magnitude after I offered to pitch in with unloading
the new cluster, and an hour of manual labor later I had unpacked a grand total of 20
blades.
The PDL retreat for being the wonderful event it is.
While it took me a long time to
understand how to leverage it well,
this work would not have been possible without the conversations I got to have there.
Andy, while unlikely to have been aware of this work until I asked him to be on my committee,
for an excellent seminar roster and a YouTube channel,
influences from which led to Garth suspecting a direct role in my pro-SQL advocacy.
Garth for building this enduring institution and for thought-provoking conversations during the odd
appearance.
Finally, the students: Rajat, Angela, Michael, Saurabh, Jack, Pratik, Juncheng, Sara, Val, Daniel,
Hojin, Daiyaan, Nj, Ziyue, Jekyeom, Sheng, Sherman, Sanjith, Wan, Sam, Theo, and Tim.
Time, a plague, and age unfortunately make exhaustive recollection difficult.

\looseness=-1
I was fortunate to infiltrate cliques beyond PDL through some strategic maneuvering.
First, networks counterparts who did not ask questions once I learned to
shake my head disappointingly while muttering BBR: Philip, Nirav, Hugo, Anup, Shawn, and Chris.
Wiselab folks who learned to tolerate me after I refused to leave CIC 2300:
Adwait, Artur, John, Akarsh, Arjun, Tianshu, Edward, Tarana, Sagar, Shruti, Mallesh, and Humza.
Tao for reliable company and for the long marches to Bao and back.
Anthony Rowe for funding coffee, 3D printers, and other toys --- may the sun shine brightly
on his solar panels.
Among the broader Cylab/CIC community: Yoshi for teaching me
Wittgenstein's takes on paradoxes, Michael Stroucken for the scholar charm that
seems to be loosely correlated with progress accelerating, Maverick for being a fellow
hueligan, Milind, Lucas, and Theo for the serendipitously formed observability cabal,
Brigette for the periodic coffee resupplies,
and Zhuo for nothing in particular.

Next, fellow MSR alumni who wandered into CMU contemporaneously and friends I made through them
--- Sandeep, PV, Suhas, Don, Chirag, Philip, Anish, Raut, and Boby Robert.
A random Walmart tennis impulse that Sandeep and PV went along on led to a vibrant tennis scene
through which I met many people and burned
many calories, although substantially offset by the Page's trips that typically followed after.
Philip and his roommates Miguel, Ioannis, and Maggie opened their backyard to my Mazda, where it
acquired many brake pads, belts,
and control arms, setting an example for other backyard-havers.
Likewise for Don and Suhas for abandoning a gaming session and appearing with a spare wheel
when one was needed.
Chintu for provisioning the consideration for provisions that were instrumental to the
job search effort,
and Philip (again) for brokering access to the aforementioned provisions.
Anup for exhaustive and transparent accounts of his job search process, replete with relevant
figures, replicating the methods wherein with a fraction of the diligence proved adequate in
securing employment.

\looseness=-1
Next, businesses, communities, and people that helped build, tune, and maintain my vehicles and
myself.
Rocky at Kraynick's bike shop for helping build my e-bike and maintain it over the years.
Harbor Freight, Rock Auto, and AliExpress for efficient, value-adding operating models.
UPMC for being a one-stop shop
for systems bottlenecks within myself, and Pitt Dental for picking up the minor residual slack.
Highland Park Tennis Club for the wonderful and free summer clinics.
Justin and TechSpark for teaching me
woodworking---not being able to use the wood CNC router there remains a minor regret.
Sandeep, again, for loaning his Mazda as a source of sensors when a notoriously tedious P0171
code appeared
in mine.
Soylent and Huel for replicable and reasonably complete nutritional formulations.
Ami for a meme that slowly acquired salience, leading to substantial personal changes around Spring
2023.

% expand page by one line
% \enlargethispage{2\baselineskip}

Finally, my roommate Jovina and my parents.
Jovina for a fascinating and entertaining character study, teaching
me patience, deep breathing, and the ability to find joy in sub-par humor and random babies passing
by.
My parents for bottling up any skepticism and befuddlement that inevitably accumulated along with
the years,
and for the moderate but durable lifestyle reforms they undertook to halt the march of
atherosclerotic cardiovascular
disease.
It is finally done, and a new chapter of gainful employment beckons.

% \vfill \noindent\emph{This work was supported in part by the PDL Consortium (Bloomberg LP,
% Everpure, Google, Jane Street, LayerZero Labs, Meta, Microsoft Research, Oracle Corporation,
% Salesforce, Uber, and Western Digital), whose interest, insights, and feedback helped shape this
% work.
% }

% {\footnotesize\noindent\emph{
%     This research was sponsored by the Los Alamos National Laboratory's Institute for Reliable
%     High Performance Information Technology (IRHPIT), the NetSketch project, and the Mochi project
%     (DOE/NNSA Contract Nos.~DE-AC52-06NA25396/394903, DE-AC52-06NA25396/394957, and
%     89233218CNA000001/394957 respectively), and the New Mexico Consortium.
%     I also thank the
%     current and former members of the Parallel Data Lab (PDL) Consortium at Carnegie Mellon
%     University (Bloomberg LP, Datadog, Everpure, Google, Jane Street, LayerZero Labs, Meta,
%     Microsoft Research, Oracle Corporation, Salesforce, Uber, and Western Digital) for their
%     interest, insights, and feedback.
%   The views and the conclusions are those of the author and do not necessarily represent the opinions of the sponsoring organizations, agencies, or the U.S. Government.
% }
%   \par
% }
